% =============================================================
% บทที่ 6 การดูแลระบบ
% ข้อมูลตรวจจาก main repo: deploy/ (deploy.sh, docker-compose.yml), .github/workflows,
% CLAUDE.md (migration caveat), backend/scripts/create-admin.js, backend/supabase (Postgres 17)
% =============================================================
\label{ch:admin}

\section{งานประจำของผู้ดูแลระบบ}

งานที่ผู้ดูแลระบบต้องทำเป็นประจำแสดงดังตารางที่~\ref{tab:routine}

\begin{table}[H]
\caption{งานประจำของผู้ดูแลระบบ}
\label{tab:routine}
\centering
\begin{tabular}{|L{7.4cm}|L{3.4cm}|C{2.4cm}|}
\hline
\textbf{งาน} & \textbf{ความถี่} & \textbf{หัวข้อ} \\ \hline
ตรวจคิวรายงานข้อมูล คำขอกิจกรรม และสถานที่จาก Google & ทุกวัน & \ref{sec:admin-reports} \\ \hline
อัปเดตดัชนีค้นหาของแชตบอต & ทุกครั้งที่แก้ไขข้อมูล & \ref{sec:admin-knowledge} \\ \hline
ตรวจผลการซิงก์อัตโนมัติและค่าใช้จ่ายของ Google & ทุกสัปดาห์ & \ref{sec:auto-sync} \\ \hline
ตรวจสถานะบริการและพื้นที่ดิสก์ & ทุกสัปดาห์ & \ref{sec:monitor} \\ \hline
ปรับรุ่นระบบ & เมื่อมีรุ่นใหม่ & \ref{sec:upgrade} \\ \hline
\end{tabular}
\end{table}

\section{การจัดการผู้ใช้และสิทธิ์}

ระบบมี 2 บทบาทในฐานข้อมูล ได้แก่ \cmd{tourist} (สมาชิก) และ \cmd{admin} (ผู้ดูแลระบบ) สมาชิกสมัครเองผ่านหน้าเว็บ (หัวข้อ~\ref{sec:register}) ส่วนการระงับ ลบ กู้คืนบัญชี ปรับพอยท์ และส่งอีเมลยืนยันใหม่ทำผ่านหน้าผู้ดูแลระบบตามหัวข้อ~\ref{sec:admin-users}

\subsection{การเพิ่มผู้ดูแลระบบ}

รันสคริปต์ตามหัวข้อ~\ref{sec:create-admin} บนเครื่องที่มีซอร์สโค้ดและไฟล์ \filename{backend/.env} ที่ชี้ไปยังฐานข้อมูลจริง ถ้าอีเมลมีบัญชีสมาชิกอยู่แล้ว สคริปต์จะเปลี่ยนบัญชีนั้นเป็นผู้ดูแลระบบ

\begin{tip}
ควรมีผู้ดูแลระบบอย่างน้อย 2 บัญชี แดชบอร์ดจะเตือนเมื่อมีเพียงบัญชีเดียว เพราะถ้าลืมรหัสผ่านของบัญชีเดียวจะต้องรันสคริปต์ใหม่เพื่อกู้
\end{tip}

\subsection{การตั้งรหัสผ่านใหม่ให้ผู้ใช้}

สมาชิกที่ลืมรหัสผ่านตั้งใหม่ได้เองตามหัวข้อ~\ref{sec:forgot-password} ผู้ดูแลระบบไม่ต้องตั้งรหัสผ่านให้ ส่วนผู้ดูแลระบบที่ลืมรหัสผ่าน ให้รันสคริปต์ \cmd{create-admin.js} ด้วยอีเมลเดิมและรหัสผ่านใหม่

\section{การปรับรุ่นระบบ}
\label{sec:upgrade}

ระบบปรับรุ่นอัตโนมัติผ่าน GitHub Actions เมื่อรวมโค้ดเข้าสาขา \cmd{main} (หัวข้อ~\ref{sec:github-actions}) หรือปรับรุ่นด้วยตนเองตามขั้นตอนต่อไปนี้

\begin{steps}
  \item ถ้ารุ่นใหม่มีไฟล์ migration ใหม่ใน \filename{backend/supabase/migrations} ให้นำไปใช้ทีละไฟล์ โดยแต่ละไฟล์รันในทรานแซกชันเดียว
\begin{terminal}
set -a; . ~/dino/deploy/backend.env; set +a
docker run --rm -v ~/dino/backend/supabase/migrations:/m postgres:17 \
  psql "$DATABASE_URL" -1 -f /m/<ชื่อไฟล์ migration>.sql
\end{terminal}
  \item รันสคริปต์ปรับรุ่น
\begin{terminal}
bash ~/dino/deploy/deploy.sh
\end{terminal}
        \result{สคริปต์แสดงข้อความ ``deploy ok'' และการ์ดเวอร์ชันระบบบนแดชบอร์ดแสดงเวอร์ชันใหม่}
  \item ตรวจสอบระบบตามหัวข้อ~\ref{sec:verify}
\end{steps}

\begin{important}
\cmd{deploy.sh} ไม่นำ migration ไปใช้ให้ และฐานข้อมูลที่ใช้งานจริงไม่มีประวัติ migration จึงห้ามใช้ \cmd{npx supabase db push} กับฐานข้อมูลนี้ หลังนำ migration ไปใช้ ให้ตรวจ policy และสิทธิ์ของตารางที่เกี่ยวข้องอีกครั้ง
\end{important}

สคริปต์ \filename{deploy/deploy.sh} ทำงานตามลำดับดังนี้

\begin{enumerate}
\item ดึงโค้ดล่าสุดของสาขา \cmd{main} จาก GitHub โดยอนุญาตเฉพาะการรวมแบบ fast-forward
\item อ่านเลขเวอร์ชันจากไฟล์ \filename{VERSION} และรหัส commit เพื่อแสดงในหน้าผู้ดูแลระบบ
\item สร้าง image ของบริการ \cmd{chatbot}, \cmd{backend} และ \cmd{caddy} ทีละบริการ เพื่อไม่ให้หน่วยความจำของเครื่องแม่ข่ายไม่พอ
\item เริ่มบริการด้วย \cmd{docker compose up -d} แล้วสั่งให้ Caddy อ่านไฟล์ตั้งค่าใหม่
\item ตรวจ \cmd{/health} ซ้ำไม่เกิน 3 นาที ถ้าสำเร็จแสดง ``deploy ok'' ถ้าไม่สำเร็จแสดงบันทึกการทำงานล่าสุดของ \cmd{chatbot} และจบด้วยสถานะผิดพลาด
\end{enumerate}

\subsection{การปรับรุ่นอัตโนมัติด้วย GitHub Actions}
\label{sec:github-actions}

โครงงานใช้ GitHub Actions ที่กำหนดไว้ในโฟลเดอร์ \filename{.github/workflows} ตามตารางที่~\ref{tab:workflows}

\begin{table}[H]
\caption{ขั้นตอนอัตโนมัติของ GitHub Actions}
\label{tab:workflows}
\centering
\begin{tabular}{|L{3.4cm}|L{4.4cm}|L{5.4cm}|}
\hline
\textbf{ไฟล์} & \textbf{เงื่อนไขที่ทำงาน} & \textbf{สิ่งที่ดำเนินการ} \\ \hline
\filename{ci-frontend.yml} & pull request หรือการรวมเข้าสาขา \cmd{main} ที่แก้ไขโฟลเดอร์ \filename{frontend} & ติดตั้งไลบรารีด้วย \cmd{npm ci} และสร้างไฟล์ด้วย \cmd{npm run build} \\ \hline
\filename{ci-chatbot.yml} & pull request หรือการรวมเข้าสาขา \cmd{main} ที่แก้ไขโฟลเดอร์ \filename{chatbot-service} & ติดตั้ง Python 3.11 และไลบรารี แล้วรัน \cmd{pytest} โดยใช้ค่าจำลองแทนคีย์จริง \\ \hline
\filename{deploy-chatbot.yml} & การรวมเข้าสาขา \cmd{main} ที่แก้ไข \filename{frontend}, \filename{backend}, \filename{chatbot-service} หรือ \filename{deploy} หรือสั่งเริ่มเอง & เชื่อมต่อเครื่องแม่ข่ายผ่าน SSH แล้วรัน \cmd{deploy.sh} \\ \hline
\end{tabular}
\end{table}

ก่อนใช้การปรับรุ่นอัตโนมัติ ต้องเพิ่มค่าลับที่ \menu{Settings > Secrets and variables > Actions} ของ repository ตามตารางที่~\ref{tab:gh-secrets}

\begin{table}[H]
\caption{ค่าลับของ GitHub Actions สำหรับการปรับรุ่นอัตโนมัติ}
\label{tab:gh-secrets}
\centering
\begin{tabular}{|L{4.2cm}|L{9cm}|}
\hline
\textbf{ชื่อค่าลับ} & \textbf{คำอธิบาย} \\ \hline
\cmd{VPS_HOST} & ชื่อโฮสต์หรือที่อยู่ของเครื่องแม่ข่าย \\ \hline
\cmd{VPS_USER} & ชื่อผู้ใช้ที่ใช้เชื่อมต่อ SSH \\ \hline
\cmd{VPS_SSH_KEY} & กุญแจส่วนตัวของคู่กุญแจที่สร้างสำหรับ CI เท่านั้น \\ \hline
\cmd{VPS_KNOWN_HOSTS} & ผลลัพธ์ของคำสั่ง \cmd{ssh-keyscan -t ed25519} ตามด้วยค่า \cmd{VPS_HOST} สำหรับยืนยันตัวตนของเครื่องแม่ข่าย \\ \hline
\end{tabular}
\end{table}

กุญแจสาธารณะของคู่กุญแจดังกล่าวต้องเพิ่มในไฟล์ \filename{~/.ssh/authorized_keys} บนเครื่องแม่ข่าย โดยกำหนดคำสั่งบังคับ (forced command) ให้รันได้เฉพาะ \cmd{deploy.sh} แม้กุญแจรั่วไหลก็เปิดเชลล์เข้าเครื่องไม่ได้ รูปแบบของบรรทัดดังกล่าวอยู่ในความคิดเห็นต้นไฟล์ \filename{deploy/deploy.sh}

\begin{warning}
ขั้นตอนตรวจสอบ (CI) และขั้นตอนปรับรุ่น (deploy) ทำงานแยกกัน การปรับรุ่นจึงไม่รอผลการทดสอบ ผู้พัฒนาต้องตรวจให้การทดสอบผ่านก่อนรวมโค้ดเข้าสาขา \cmd{main} และ backend ยังไม่มีขั้นตอน CI ต้องรัน \cmd{npm test} เองก่อนรวมโค้ด
\end{warning}

\section{การตรวจสถานะและบันทึกการทำงาน}
\label{sec:monitor}

ตรวจสถานะของระบบด้วยคำสั่งในตารางที่~\ref{tab:monitor} ที่โฟลเดอร์ \filename{~/dino/deploy} นอกจากนี้การ์ดเวอร์ชันระบบบนแดชบอร์ดแสดงเวอร์ชันและรหัส commit ของทุกบริการ ใช้ตรวจว่าปรับรุ่นสำเร็จหรือไม่

\begin{table}[H]
\caption{การตรวจสถานะของระบบ}
\label{tab:monitor}
\centering
\begin{tabular}{|L{3.2cm}|L{5.4cm}|L{4.6cm}|}
\hline
\textbf{สิ่งที่ตรวจ} & \textbf{คำสั่ง} & \textbf{ผลที่ปกติ} \\ \hline
สถานะบริการ & \cmd{docker compose ps} & \cmd{caddy}, \cmd{backend} และ \cmd{chatbot} มีสถานะ Up \\ \hline
backend และฐานข้อมูล & \cmd{curl https://<CHAT_DOMAIN>/api/health} & \cmd{{"status":"ok"}} \\ \hline
AI Service & \cmd{curl https://<CHAT_DOMAIN>/health} & \cmd{{"status":"ok"}} \\ \hline
บันทึกการทำงาน & \cmd{docker compose logs --tail 100 backend} และ \cmd{chatbot} & ไม่มีข้อความ Error ซ้ำ ๆ \\ \hline
พื้นที่ดิสก์ & \cmd{df -h} และ \cmd{docker system df} & พื้นที่ว่างมากกว่าร้อยละ 20 \\ \hline
\end{tabular}
\end{table}

\begin{tip}
image รุ่นเก่าที่เหลือจากการปรับรุ่นกินพื้นที่ดิสก์ ลบได้ด้วยคำสั่ง \cmd{docker image prune -f}
\end{tip}

\section{ข้อควรระวังด้านความปลอดภัย}

\begin{enumerate}
\item ไฟล์ \filename{.env}, \filename{backend.env} และ \filename{chatbot.env} ต้องไม่ถูกบันทึกเข้า Git และควรจำกัดสิทธิ์ด้วยคำสั่ง \cmd{chmod 600}
\item \cmd{SUPABASE_SERVICE_ROLE_KEY} ข้ามการตรวจสิทธิ์ของฐานข้อมูลได้ทั้งหมด ใช้เฉพาะในฝั่งเครื่องแม่ข่าย ห้ามใส่ในตัวแปรของ frontend
\item จำกัดคีย์ Google Maps ของ frontend ด้วย HTTP referrer ของเว็บไซต์ และตั้งการแจ้งเตือนงบประมาณใน Google Cloud สำหรับคีย์ Places ของ backend
\item เปิดเฉพาะพอร์ต 80, 443 และ SSH บนเครื่องแม่ข่าย พอร์ตของ backend และ AI Service ต้องเข้าถึงได้ภายในเครือข่าย Docker เท่านั้น
\item ใช้รหัสผ่านที่คาดเดายากสำหรับบัญชีผู้ดูแลระบบ และมีผู้ดูแลระบบอย่างน้อย 2 บัญชี
\item กุญแจ SSH ของ CI ต้องมีคำสั่งบังคับให้รันได้เฉพาะ \cmd{deploy.sh}
\item ถ้าใช้ Vertex AI ให้ใช้ service account ที่มีเฉพาะสิทธิ์ \cmd{roles/aiplatform.user} และกำหนดสิทธิ์ไฟล์กุญแจเป็น 0400
\item ถ้าสงสัยว่าคีย์รั่วไหล ให้สร้างคีย์ใหม่ที่ผู้ให้บริการ แก้ไฟล์ตั้งค่า แล้วรัน \cmd{docker compose up -d} และสั่ง \btn{ออกจากระบบทุกอุปกรณ์} ให้บัญชีที่เกี่ยวข้อง
\end{enumerate}
